\section{Introduction}\label{sec:introduction}

Single-site heat-bath dynamics is a natural way to sample an interacting
spin system: one spin at a time is resampled from its conditional Gibbs
law. In the Sherrington--Kirkpatrick model~\cite{SK1975}, every pair of
spins interacts through a Gaussian coupling. At low temperature, a
configuration chosen adversarially can be much harder to escape than a
configuration drawn from equilibrium. Indeed, at sufficiently low
temperature, gapped configurations force an exponential worst-start
mixing-time lower bound~\cite[Theorem~1.1]{Sellke2025}.

We prove an upper bound throughout the fixed low-temperature phase for a
different question. Sample the disorder, sample one configuration from
its Gibbs measure, and then fix that configuration as the initial state.
How long does its transition law take to approach equilibrium? The
answer is at most $\exp(n^{1-c})$ with probability tending to one, for
an explicit positive $c$. This formulation keeps the randomness of the
initial state outside the total-variation distance.

\subsection{Model and main theorem}

Let $n\ge2$, and let $J_{ij}$, $1\le i<j\le n$, be independent
standard Gaussian variables. For $\sigma\in\{-1,1\}^n$, define
\begin{equation*}\tag{H}\label{eq:intro-model}
 H_n^J(\sigma)=\frac1{\sqrt n}\sum_{i<j}J_{ij}\sigma_i\sigma_j,
 \qquad
 \pi_{n,\beta}^J(\sigma)
 =\frac{\exp(\beta H_n^J(\sigma))}{Z_{n,\beta}^J}.
\end{equation*}
The external field is zero. Each site has an independent rate-one clock.
When the clock at site $i$ rings, the new spin equals $s\in\{-1,1\}$
with conditional probability
\[
 \frac{\exp(\beta s h_i(\sigma))}{2\cosh(\beta h_i(\sigma))},
 \qquad
 h_i(\sigma)=\frac1{\sqrt n}\sum_{j\ne i}J_{ij}\sigma_j,
\]
where $J_{ji}=J_{ij}$ and $J_{ii}=0$. Such an update may leave the
configuration unchanged. Denote the resulting continuous-time
transition law from a fixed $\sigma$ by $P_t^{J,\beta}(\sigma,\cdot)$,
and set
\begin{equation*}\tag{T}\label{eq:intro-state-mixing-time}
 T_{n,\beta}^J(\sigma)
 =\inf\left\{t\ge0:
 \left\|P_t^{J,\beta}(\sigma,\cdot)-\pi_{n,\beta}^J\right\|_{\mathrm{TV}}
 \le\frac14\right\}.
\end{equation*}

\begin{theorem}[Mixing from a typical Gibbs state]\label{thm:main}
For every fixed $\beta>1$, there is $c_\beta\in(0,1)$ such that
\[
 \mathbb P_{J,\,\sigma_0\sim\pi_{n,\beta}^J}
 \left\{T_{n,\beta}^J(\sigma_0)
       \le\exp\bigl(n^{1-c_\beta}\bigr)\right\}
 \longrightarrow1
 \qquad\text{as }n\longrightarrow\infty.
\]
The exponent is independent of $n$ and of the realized disorder.
The proof permits $c_\beta=1/40000000$ for every fixed $\beta>1$.
\end{theorem}

Equivalently, in probability over $J$, the Gibbs mass of states
satisfying the displayed bound tends to one. The state $\sigma_0$ is
realized before the transition law and its distance from equilibrium
are formed. If the initial state were averaged into that transition law,
stationarity would instead make the distance zero at every time.
The exponent above is not optimized, and the dimension at which the
estimates become effective may depend on the fixed $\beta$. All times
in this paper use the continuous-time convention just specified.

\begin{corollary}[Two-sided typical-start bounds]\label{cor:typical-two-sided}
For every fixed $\beta>1$, in the model and continuous-time convention
above,
\[
 \mathbb P_{J,\,\sigma_0\sim\pi_{n,\beta}^J}
 \left\{
 \begin{aligned}
  \exp\bigl(n^{1/10000}\bigr)&<T_{n,\beta}^J(\sigma_0)\\
  &\le\exp\bigl(n^{1-1/40000000}\bigr)
 \end{aligned}
 \right\}\longrightarrow1
 \qquad\text{as }n\longrightarrow\infty.
\]
The Gibbs state $\sigma_0$ is held fixed before the total variation in
\eqref{eq:intro-state-mixing-time} is evaluated.
\end{corollary}

\begin{proof}
The continuous-time part of \cite[Theorem~1.1]{OpenAISKBarriers2026}
gives total-variation distance greater than $1/4$ at time
$\exp(n^{1/10000})$ with joint probability tending to one. Contraction
of total variation and continuity of the finite-state semigroup then
give the strict lower bound on $T_{n,\beta}^J(\sigma_0)$.
The upper bound is Theorem~\ref{thm:main} with
$c_\beta=1/40000000$. Both events use the same disorder/Gibbs-start law
and rate-one-per-site clocks, so their intersection has probability
tending to one by the union bound.
\end{proof}

\subsection{Earlier mixing results and proof ingredients}

At high temperature, functional inequalities have yielded polynomial
mixing bounds from every starting state. Eldan, Koehler and Zeitouni
proved such a bound for $\beta<1/4$ using a spectral condition on the
interaction matrix~\cite[Section~6]{EKZ2022}. More recently, Wang
proved optimal-order Glauber mixing for every fixed $\beta<1/2$
\cite[Theorem~1.1]{Wang2026}, and Boban, Li and Oveis Gharan
extended polynomial mixing in zero field to
$\beta<1/2+\varepsilon_0$ for an absolute $\varepsilon_0>0$
\cite[Corollary~1.6]{BLO2026}. These results concern worst-start
mixing in a high-temperature range. In the low-temperature setting,
Sellke explicitly distinguished the worst-start obstruction from the
possibility of faster mixing after Gibbs-distributed
initialization~\cite[discussion after Theorem~1.1]{Sellke2025}.
Theorem~\ref{thm:main} addresses this latter question throughout the
range of fixed $\beta>1$.

The equilibrium input that drives the proof is stationarity along an
entire interval of scalar overlaps. The variational analysis of the
Parisi functional developed by Chen~\cite[Theorem~2 and Proposition~1]{Chen2017} and by Jagannath
and Tobasco~\cite[Proposition~1.1 and Corollaries~3.6,~3.10]{JT2017}
provides the scalar stationarity conditions on the support of the
minimizing measure. Zhou established full interval support immediately
above the transition~\cite[Theorem~1]{Zhou2026}. For the whole range
of fixed $\beta>1$, we use the full interval conclusion of
Lopatto~\cite[Theorem~1.1]{Lopatto2026}, specifically version~3.
Full support lets us apply the stationarity identity at every overlap
between zero and the top support point. Infinite support alone would
not supply this input.

Two other equilibrium results enter with different purposes. Guerra's
finite-volume interpolation bound~\cite[Theorem~3]{Guerra2003}
controls the expected log partition function from above. The spatial
regularity and step-approximation stability of the scalar Parisi
equation~\cite[Propositions~1--2]{AC2015} justify the diffusion and its
calculus. Section~\ref{sec:scalar} states these inputs in the exact
normalization used here and derives the additional estimates beyond
the support endpoint.

The proposal construction is related to established Gaussian-query
methods. Bolthausen's iterative construction of TAP solutions uses
successive Gaussian conditioning and orthogonalization
\cite[Sections~2--3]{Bolthausen2014}. Montanari's incremental approximate
message passing algorithm uses a Parisi diffusion and normalized
magnetization increments to optimize the SK Hamiltonian
\cite[Sections~2--3]{Montanari2021}. Our proposal uses exact
finite-dimensional Gaussian conditioning to establish quantitative
entropy and likelihood comparisons with the Gibbs joint law.

The comparison in the reverse direction uses transport under a
strongly log-concave reference measure
\cite[Corollary~3.10]{BKM2005}, followed by Gaussian smoothing
\cite[Lemma~11]{LugosiNeu2022}. We give the necessary conditional
proofs in Section~\ref{sec:reverse}. Finally, the conversion from
flow to mixing uses the coarea proof of the reversible Cheeger bound
\cite[Theorem~3.5]{LawlerSokal1988} and entropy dissipation
\cite[Lemmas~2.5 and~2.7]{DiaconisSaloffCoste1996}. Our finite-chain
argument retains a set of large stationary mass before making this
conversion, so that it only requires information about cuts separating
appreciable masses.

\subsection{How the proof reaches a typical fixed state}

We first construct a random spin configuration correlated with the
Gaussian interaction matrix. Its coordinates follow a scalar
magnetization martingale. The martingale continues beyond the Parisi
support until the average coordinate uncertainty is small. During this continuation,
the matrix-query noise is mixed with an independent Brownian noise.
An exact scalar identity balances the entropy paid through the matrix
channel against the energy gained by the proposal. Sections
\ref{sec:scalar} and~\ref{sec:proposal} turn this identity into two
estimates: sublinear forward relative entropy, and an upper likelihood
bound outside an exponentially unlikely event.

Forward relative entropy by itself does not guarantee that a proposal
reaches a prescribed set of positive Gibbs mass. To obtain that
property, Section~\ref{sec:reverse} conditions the Gibbs joint law on
the spin. The remaining disorder law is strongly log-concave. Smoothing
the disorder then reverses the entropy comparison and gives a lower
bound on the proposal probability of every set whose Gibbs mass is
appreciable, for the realized target disorder and smoothing base.

Section~\ref{sec:connections} joins two proposals generated independently
conditional on that base. Longer, fine-horizon proposals supply the
endpoint coverage; shorter, coarse-horizon proposals have the rarer
failures needed along the connecting sequence. Near-pinning of the
magnetizations joins these two scales. Rotating the Gaussian inputs of
the coarse proposals produces small consecutive Hamming distances.
If a cut separating appreciable Gibbs masses had very small
stationary flow, its enlarged vertex boundary would have tiny Gibbs
mass. The likelihood bound makes the intermediate proposals avoid that
boundary, while coverage gives a positive chance that the two endpoints
lie on opposite sides. These conclusions contradict the elementary
fact that a short-step sequence crossing a cut must approach its
boundary; Figure~\ref{fig:connection} depicts this last step.

It remains to turn the resulting flow estimate into a mixing statement.
Section~\ref{sec:mixing} proves a finite-chain lemma that removes less
than a prescribed small stationary mass and obtains a Poincar\'e
inequality on the remainder. Entropy dissipation and stationarity then
bound the average of the \emph{fixed-state} total-variation distances.
The argument does not require a trajectory to stay in the retained set.
This is the final reason that rare slow starting states do not prevent
the conclusion of Theorem~\ref{thm:main}.
